\documentclass{article}

\usepackage{tabularx}
\usepackage{booktabs}

\title{Problem Statement and Goals\\\progname}

\author{\authname}

\date{}

\input{../Comments.text}
\input{../Common.text}

\begin{document}

\maketitle

\begin{table}[hp]
\caption{Revision History} \label{TblRevisionHistory}
\begin{tabularx}{\textwidth}{llX}
\toprule
\textbf{Date} & \textbf{Developer(s)} & \textbf{Change}\\
\midrule
Date1 & Name(s) & Description of changes\\
Date2 & Name(s) & Description of changes\\
... & ... & ...\\
\bottomrule
\end{tabularx}
\end{table}

\section{Problem Statement}

\wss{You should check your problem statement with the
\href{https://smiths.github.io/capTemplate/Checklists/ProbState-Checklist.pdf}
{problem statement checklist}}. 

\wss{You can change the section headings, as long as you include the required
information.}

\subsection{Problem}

\wss{What problem are you trying to solve?  Why is it important? Do not talk
about how you will solve the problem, only what the problem is.  AI is rarely
part of the problem.  The problem is ``what'' you want to do, not ``how'' you
want to do it.  The problem statement should be written in a way that is
understandable to a non-technical audience.}

\subsection{Inputs and Outputs}

\wss{Characterize the problem in terms of ``high level'' inputs and outputs.  
Use abstraction so that you can avoid details.}

\subsection{Stakeholders}

\subsection{Environment}

\subsubsection{Software}
The system will be a web-based application accessed through the existing TeleHealth
platform, and will be supported on all current web browsers. The platform itself will
be built using modern web frameworks such as React, and languages like Python.
Modules will follow common software design patterns and object-oriented programming
principles. Additional technologies and libraries will also be used for purposes such
as cloud services in order to host the platform, data storage, and authentication.

\subsubsection{Hardware}
Clinicians will access the annotation interface from a desktop view,
as reviewing and editing transcripts alongside video playback requires a larger
screen, as well as a keyboard and mouse for precise editing. Audio output, such as
speakers or headphones, will also be required to review speech and videos
from recorded assessments. Patients and caregivers will complete assessment
sessions through TeleHealth's existing web-based interface, using a device with a 
webcam and microphone. Support for assessments and annotations on mobile phones
or tablets will be a stretch goal but may likely be out of scope for this project.

\section{Goals}

\section{Stretch Goals}

\wss{Goals you hope to get to, but not necessary.  They are also helpful if the
scope of the original project needs to be expanded.}

\section{Custom Documents (CDs)}

\wss{For CAS 741: State whether the project is a research project. This
designation, with the approval (or request) of the instructor, can be modified
over the course of the term.}

\wss{For SE Capstone: List your custom documents (CDs).  Potential CDs include
usability testing, code walkthroughs, user documentation, formal proof,
GenderMag personas, Design Thinking, etc.  (The full list is on the course
outline and in Lecture 02.) Normally the number of CDs will be two (one per
term).  Approval of the CDs will be part of the discussion with the instructor
for approving the project. The CDs, with the approval (or request) of the
instructor, can be modified over the course of the term.}

\newpage{}

\section*{Appendix --- Reflection}

\wss{Not required for CAS 741}

\input{../Reflection.text}

\begin{enumerate}
    \item What went well while writing this deliverable? 
    \item What pain points did you experience during this deliverable, and how
    did you resolve them?
    \item How did you and your team adjust the scope of your goals to ensure
    they are suitable for a Capstone project (not overly ambitious but also of
    appropriate complexity for a senior design project)?
\end{enumerate}  

\subsection*{Catherine}
Our team was given a wide breadth of resources and knowledge from both our supervisor
and additional professors also working on the MERS project, so we were able to define
our problem and goals pretty well before starting this deliverable. We also
planned multiple meetings in order to discuss each task item, brainstorm various ideas/notes, 
and ensure that everybody was on the same page before we broke off to
individually tackle our assigned parts. Having both the project knowledge and pre-work
discussions helped heavily in guiding us through our deliverable, so we were able
to easily communicate our individual thoughts with each other and work through the
deliverable efficiently.

One of the pain points I believe we experienced was understanding the nitty-gritty
details of our project, given how it touches several moving parts, from the SAFA
system to the TeleHealth platform. There were many questions we had about the project,
and we often had to wait for responses, gather information from lengthy papers, and
make our own assumptions in the meantime. What helped a lot was, again, our frequent
and productive meetings, where we brainstormed and worked on understanding the project
together. Once we had a better grasp of the project, another pain point was then figuring 
out our scope; e.g what can we realistically finish and do we need to scope up or 
scope down? To work through this, we outlined our main goals and separated those from our
stretch goals, and plan on using our POC as a checkpoint since it validates
our most concerning risks. 

To keep our scope as realistic but properly complex as possible, we made sure that as
a team, we fully and thoroughly understood the project and what was being asked of us.
This included understanding what parts of the MERS system were done, what could
be improved, and what needed to still be implemented or connected. By separating those 
into individual goals as previously mentioned, we were able to visualize and
better understand the scope of our project and what could be doable within the next
two terms. We wanted to ensure that we weren't underestimating any technical work
involved, so actually working on this deliverable also brought up multiple 
useful questions that we asked ourselves; such as what are the risks that come with 
your project and are we able to prevent them? Essentially, visualizing our project
in a more in-depth manner, and breaking the tasks that we'd eventually
work on developing into "modules" helped us a lot in determining what would be
achievable this year.

\end{document}